%% ma: L12-B13-0034
%% lop: 12
%% chuong: 4
%% bai: 13
%% dang: 12.13.17
%% loai: TLN
%% muc: VDC
%% dap_an: "117"
%% nguon: "(NBV)-ĐỀ SỐ 16-MỖI NGÀY 1 ĐỀ THI 2025.pdf (D16, MỖI NGÀY 1 ĐỀ - Đề số 16), Phần III Câu 5"
%% kien_thuc: "Bài 12 (tích phân), Bài 13 (thể tích khối tròn xoay)"
%% trang_thai: cho-duyet
%% ly_do: "Dạng mới đề xuất 12.13.17 chờ giáo viên duyệt"
%% ghi_chu: "Đáp án gốc 117 đúng (V = 74π/3 + 9√2π ≈ 117,48). Hình dựng lại đúng tỉ lệ (OM = 20, bán kính cầu 2, MC = 1); điểm I là giao điểm của OM với đường tròn (OI = 2)."
\begin{ex}
Họa sĩ thiết kế một micro có dạng khối tròn xoay, mặt cắt đứng chứa trục của khối tròn xoay có dạng như hình sau, trong đó $OA=OB=OI=2$ cm, $MC=MD=1$ cm, đường thẳng $OM$ là đường trung trực của đoạn thẳng $CD$, $OM=20$ cm, $\widehat{AOB}=90^\circ$.
\begin{center}
\begin{tikzpicture}[scale=.62]
  \coordinate (O) at (0,0);\coordinate (A) at ({sqrt(2)},{sqrt(2)});\coordinate (B) at ({sqrt(2)},{-sqrt(2)});
  \coordinate (C) at (20,1);\coordinate (D) at (20,-1);\coordinate (M) at (20,0);\coordinate (I) at (2,0);
  \draw (A) arc[start angle=45,end angle=315,radius=2] -- (D) -- (C) -- cycle;
  \draw[khuat] (O)--(A) (O)--(B) (O)--(M) (A)--(B);
  \draw[khuat] (A) arc[start angle=45,end angle=-45,radius=2];
  \foreach \p/\d in {O/left,A/above,B/below,C/above right,D/below right,M/right,I/above right}
    \fill (\p) circle (1.2pt) node[\d,font=\footnotesize] {$\p$};
\end{tikzpicture}
\end{center}
Thể tích của micro này là bao nhiêu $\mathrm{cm}^3$? (Làm tròn kết quả đến hàng đơn vị)
\shortans{117}
\loigiai{Đặt hệ trục $Oxy$ với $O$ là gốc, tia $Ox$ là $OM$, đơn vị $1$ cm. Khi đó $A(\sqrt2;\sqrt2)$, $C(20;1)$. Đường sinh gồm cung tròn $y=\sqrt{4-x^2}$ với $x\in[-2;\sqrt2]$ và đoạn $AC$.

Đường thẳng $AC$: $y=\sqrt2+\dfrac{1-\sqrt2}{20-\sqrt2}\left(x-\sqrt2\right)$.

\[V=\pi\int_{-2}^{\sqrt2}(4-x^2)\,\mathrm dx+\pi\int_{\sqrt2}^{20}\left[\sqrt2+\dfrac{1-\sqrt2}{20-\sqrt2}\left(x-\sqrt2\right)\right]^2\mathrm dx\]
\[V=\dfrac{(16+10\sqrt2)\pi}3+\dfrac{(58+17\sqrt2)\pi}3=\dfrac{74\pi}3+9\sqrt2\pi\approx117{,}48.\]

Làm tròn: $117\ \mathrm{cm}^3$.}
\end{ex}
